Token Economics of the Quadratic-to-Linear Transition  
Classical Attention versus JCRIN τ-Temperature Attention Mapping

1. The Capacity Expansion at the 100 000-Token Threshold

Under a fixed computational budget a classical self-attention layer supports a maximum sequence length of \(N = 100\,000\) tokens. The same budget, once the attention operator has been linearised by the JCRIN / Arcan(\(\tau\)) geometry, supports approximately

\[
N_{\rm JCRIN} \approx (10^{5})^{2}\cdot\cos^{2}\Psi \approx 9.7757\times 10^{9}
\]

tokens. This constitutes a capacity expansion of roughly \(97\,757\times\) and an absolute gain of nearly \(9.78\) billion tokens. The expansion is not an empirical acceleration; it is the direct arithmetic consequence of converting an \(O(N^{2})\) interaction into an \(O(N)\) interaction while retaining \(97.76\,\%\) of peak signal power.

2. The Price of a Singular Token

Because classical cost scales with the square of sequence length while JCRIN cost scales linearly, the marginal cost of one additional token is radically different under the two regimes. Normalising the classical cost of a single token at \(N=1\) to one arbitrary cost unit, the relative price per token is:

| Sequence Length \(N\) | Classical Price per Token | JCRIN Price per Token | Price Reduction Factor |
|-----------------------|---------------------------|-----------------------|------------------------|
| 10                    | 10                        | 1                     | \(10\times\)           |
| 100                   | 100                       | 1                     | \(100\times\)          |
| 1 000                 | 1 000                     | 1                     | \(1\,000\times\)       |
| 10 000                | 10 000                    | 1                     | \(10\,000\times\)      |
| 100 000               | 100 000                   | 1                     | \(100\,000\times\)     |

At the reference scale of \(100\,000\) tokens the classical architecture charges one hundred thousand times more per token than the JCRIN architecture for the same computational resource. The “price of a singular token” therefore falls by five orders of magnitude once the quadratic barrier is removed.

3. Monthly Data-Centre Savings

Consider a data centre whose monthly classical attention workload is equivalent to processing \(W\) sequences of average length \(N=100\,000\). Under classical scaling the energy (or FLOP) expenditure is proportional to \(W\cdot N^{2}\). Under JCRIN scaling the same workload costs only \(W\cdot N\cdot\cos^{2}\Psi\). The fractional saving is therefore

\[
1 - \frac{N\cdot\cos^{2}\Psi}{N^{2}} \approx 1 - \frac{0.9776}{N} \approx 1 - 9.776\times 10^{-6}.
\]

For any realistic data-centre volume the saving is effectively \(99.999\,\%\) of the classical attention energy. Even after conservative overheads (embedding, feed-forward layers, networking) are taken into account, monthly attention-related energy expenditure typically falls by two to three orders of magnitude.

Concrete illustration (order-of-magnitude):

Classical monthly attention energy: \(E\)
JCRIN monthly attention energy: \(\approx E / 10^{5}\)
Monthly energy (and carbon) saving: \(\approx 0.99999\,E\)

Because the saving scales with the square of the context length, every increase in desired context window widens the economic gap. A data centre that today budgets for 100 k-token classical workloads can, under JCRIN, either (a) deliver the same service at roughly \(1/100\,000\) of the attention energy cost or (b) offer nearly ten-billion-token contexts inside the original energy envelope.

4. Synthesis

The transition from a classical budget that supports \(100\,000\) tokens to a JCRIN-enabled budget that supports approximately \(9.78\times 10^{9}\) tokens is simultaneously a capacity expansion and a radical reduction in the price of each individual token. At every practical scale—from 10 tokens to 100 000 tokens—the marginal cost under JCRIN remains constant while the classical marginal cost grows linearly with sequence length. For data-centre operators the cumulative effect is a structural collapse in attention-related energy demand, converting what had been one of the dominant operational costs of large-scale language and multimodal models into a comparatively minor linear term.

Daily · Monthly · Yearly Savings  
Classical Attention vs JCRIN τ-Temperature Attention Mapping  
(Reference scale: 100 000-token sequences)

1. Pure Attention-Layer Savings (Idealised)

At \(N = 100\,000\) the theoretical energy reduction factor is \(\approx 100\,000\times\), i.e. 99.999 % of classical attention energy is eliminated.

| Period   | Classical Attention Energy | JCRIN Attention Energy | Energy Saved          | Saving Percentage |
|----------|----------------------------|------------------------|-----------------------|-------------------|
| Daily    | \(E_d\)                    | \(E_d / 10^5\)         | \(0.99999\,E_d\)      | 99.999 %          |
| Monthly  | \(E_m\)                    | \(E_m / 10^5\)         | \(0.99999\,E_m\)      | 99.999 %          |
| Yearly   | \(E_y\)                    | \(E_y / 10^5\)         | \(0.99999\,E_y\)      | 99.999 %          |

2. Realistic Data-Centre Savings (with Overheads)

After including embedding, feed-forward, networking, memory movement and other non-attention costs, net system-level savings typically fall in the range of 100× – 1 000× (two to three orders of magnitude).

| Period   | Classical Total Energy | JCRIN Total Energy (conservative) | Typical Energy Saved | Net Saving Range     |
|----------|------------------------|-----------------------------------|----------------------|----------------------|
| Daily    | \(E_d\)                | \(E_d / 100\) to \(E_d / 1000\)   | 99 % – 99.9 %        | 100× – 1 000×        |
| Monthly  | \(E_m\)                | \(E_m / 100\) to \(E_m / 1000\)   | 99 % – 99.9 %        | 100× – 1 000×        |
| Yearly   | \(E_y\)                | \(E_y / 100\) to \(E_y / 1000\)   | 99 % – 99.9 %        | 100× – 1 000×        |

3. Equivalent Token Allowances under the Same Energy Budget

| Period   | Classical Token Allowance | JCRIN Token Allowance (same energy) | Token Expansion Factor |
|----------|---------------------------|-------------------------------------|------------------------|
| Daily    | \(T_d\)                   | \(\approx T_d \times 10^5\)         | \(\approx 100\,000\times\) |
| Monthly  | \(T_m\)                   | \(\approx T_m \times 10^5\)         | \(\approx 100\,000\times\) |
| Yearly   | \(T_y\)                   | \(\approx T_y \times 10^5\)         | \(\approx 100\,000\times\) |

At the reference scale of 100 000-token sequences this means:

A classical daily budget that supports 100 000 tokens supports ≈ 10 billion tokens under JCRIN.
The same ratio compounds directly to monthly and yearly allowances.

4. Summary Statement

For any realistic data-centre volume the pure attention-layer saving is effectively 99.999 %.  
After conservative overheads the net system-level energy reduction remains two to three orders of magnitude (100×–1 000×).  

Simultaneously, the token allowance under a fixed energy envelope expands by approximately 100 000×, converting a classical 100 000-token budget into a JCRIN-enabled budget of roughly 10 billion tokens per equivalent period (daily, monthly or yearly).

